\slajdid{10-s021}
\begin{frame}[label=10-s021]{Strujni kapacitet}
\gusttekst\setlength{\abovedisplayskip}{3pt}\setlength{\belowdisplayskip}{3pt}\setlength{\abovedisplayshortskip}{2pt}\setlength{\belowdisplayshortskip}{2pt}\begin{columns}[c,onlytextwidth]\column{.26\textwidth}\centering\begin{tikzpicture}[x=.52cm,y=.7cm,font=\sitneoznake,>=Latex]\ctikzset{bipoles/length=.65cm}\draw(0,-.5)--(0,.5);\draw(-.6,0)--(0,0);\draw(0,.32)--(.7,.7)--(.7,1.3);\draw[->](0,-.32)--(.7,-.7);\draw(.7,-.7)--(.7,-1.3)node[ground]{};\node[right]at(.25,0){$T_1$};\draw(-2.3,0)node[ocirc]{}to[R,l=$R_B$](-.6,0);\draw(.7,1.3)to[R,l_=$R_C$](.7,3);\draw(.4,3)--(1,3);\node[above]at(.7,3){$V_{CC}$};\draw(.7,1.3)--(2,1.3)node[ocirc]{};\node[left]at(-2.3,0){$V_I$};\node[right]at(2,1.3){$V_O$};\draw[->](-2.2,1)--(-1.8,1)node[midway,above]{$I_I$};\draw[->](-.6,.4)--(-.15,.4)node[midway,above]{$I_{B1}$};\draw[->](-1.6,-.35)--(-1.1,-.35)node[midway,below]{$I_{RB}$};\draw[->](.35,2.55)--(.35,1.95)node[midway,left]{$I_{RC}$};\draw[->](2.1,1.6)--(1.5,1.6);\node at(2.5,2.1){$I_O$};\draw[->](.95,1.2)--(.95,.65)node[midway,right]{$I_{C1}$};\end{tikzpicture}


\column{.72\textwidth}
Da bi odredili $I_{IL}$ smatraćemo da se na ulazu nalazi napon od $V_{OL}$ do $V_{IL}$. U tom slučaju tranzistor T1 je zakočen, kao što smo videli, pa je njegova bazna struja $I_B=0$, odnosno i ulazna struja jednaka nuli. Znači $I_{IL}=0$.
\par\medskip
Da bi odredili $I_{IH}$ smatraćemo da se na ulazu nalazi napon od $V_{IH}$ do $V_{OH}$. U tom slučaju tranzistor T1 je u zasićenu, kao što smo videli, pa je njegova bazna struja $I_B$, što je isto što i ulazna struja $I_I$, po ulaznoj konturi
\[I_B=I_I=\frac{V_I-V_{BE1}}{R_B}=\frac{V_I-V_{BES}}{R_B}\]
Očigledno je ona maksimalna kada je i maksimalan ulazni napon $V_{I\max}=V_{OH}=V_{CC}$ pa je
\[I_{Imax}=I_{IH}=\frac{V_{CC}-V_{BES}}{R_B}\]
\end{columns}
\end{frame}
